% ============================================================
%  Projective Dynamic Logo — Popular science for general audiences
%  Cédric Laubscher — 2026
%  Compile with XeLaTeX or LuaLaTeX (recommended)
%  or pdfLaTeX (replace fontspec with inputenc/fontenc)
% ============================================================

\documentclass[11pt, a4paper, openright]{book}
\usepackage{hyphenat}
\usepackage{hyphenat}
% --- Encoding and language ---
\ifdefined\XeTeXversion
\usepackage{ifxetex,ifluatex}
\ifxetex
\usepackage{fontspec}
  \setmainfont{Latin Modern}           % replace with EB Garamond if unavailable
  \setsansfont{Source Sans Pro}
  \setmonofont{JuliaMono}[Scale=0.85]
\else
  \usepackage[T1]{fontenc}
  \usepackage[utf8]{inputenc}
  \usepackage{lmodern}
\fi

  \usepackage[utf8]{inputenc}
  \usepackage[T1]{fontenc}

\let\showhyphens\relax
\usepackage[english]{babel}
\usepackage{csquotes}

% --- Page layout ---
\usepackage{xcolor}
\usepackage[
  top=2.8cm, bottom=3cm,
  left=2.8cm, right=2.8cm,
  headheight=15pt
]{geometry}
\usepackage{emptypage}
\usepackage{setspace}
\onehalfspacing

% --- Typography ---
%\setmainfont{EB Garamond}           % replace with Latin Modern if unavailable
\usepackage[sfdefault]{sourcesanspro}
%\setmonofont{JuliaMono}[Scale=0.85]

% --- Mathematics ---
\usepackage{amsmath, amssymb, amsthm}
\usepackage{mathtools}

% --- Headers and footers ---
\usepackage{fancyhdr}
\pagestyle{fancy}
\fancyhf{}
\fancyhead[LE]{\small\itshape\leftmark}
\fancyhead[RO]{\small\itshape\rightmark}
\fancyfoot[C]{\small\thepage}
\renewcommand{\headrulewidth}{0.4pt}

% --- Chapter titles ---
\usepackage{titlesec}
\titleformat{\chapter}[display]
  {\normalfont\large\normalfont\centering}
  {\chaptertitlename\ \thechapter}
  {11pt}
  {\huge\centering}
\titlespacing*{\chapter}{0pt}{0pt}{30pt}

% --- Table of contents ---
\usepackage{tocloft}
\renewcommand{\cfttoctitlefont}{\normalfont\huge}
\renewcommand{\cftchapfont}{\bfseries}
\renewcommand{\cftsecfont}{\normalfont}
\setlength{\cftbeforechapskip}{6pt}


% --- Links and PDF ---
\usepackage[
  colorlinks=true,
  linkcolor=black,
  urlcolor=blue!60!black,
  citecolor=black,
  pdftitle={Whoever We May Be — The Projective Dynamic Logo},
  pdfauthor={Cédric Laubscher},
  pdfsubject={Scientific and philosophical popular science},
  pdflang={en}
]{hyperref}

% --- Epigraphs ---
\usepackage{epigraph}
\setlength{\epigraphwidth}{0.78\textwidth}
\renewcommand{\epigraphflush}{center}
\renewcommand{\epigraphrule}{0.4pt}

% --- Custom lists ---
\usepackage{enumitem}
\setlist[description]{leftmargin=2em, labelindent=0pt, font=\normalfont}

% --- Paragraph spacing ---
\setlength{\parindent}{1.4em}
\setlength{\parskip}{0pt}

% --- Custom commands ---
\newcommand{\ldp}{\textsc{pdl}}
\newcommand{\eg}{\textit{e.g.}\ }
\newcommand{\ie}{\textit{i.e.}\ }

% ============================================================
\begin{document}

% --- Title page ---
\begin{titlepage}
  \centering
  \vspace*{3cm}
  {\fontsize{11}{13}\selectfont Cédric Laubscher}\\[1.2cm]
  {\fontsize{30}{36}\selectfont\normalfont Whoever We May Be}\\[0.6cm]
  {\fontsize{16}{20}\selectfont\itshape
    An Introduction to the Projective Dynamic Logo}\\[0.4cm]
  {\fontsize{13}{16}\selectfont\itshape
    From the Electron to the Human Condition:\\
    the logic of coherence that governs everything}\\[3cm]
  \rule{0.45\textwidth}{0.4pt}\\[0.8cm]
  {\small Independent Researcher, Switzerland}\\[0.3cm]
  {\small March 2026}\\[1.5cm]
  \vfill
  {\footnotesize
    All technical publications of the PDL are available\\
    in open access on \href{https://zenodo.org}{Zenodo}
    under the name of Cédric Laubscher.}
\end{titlepage}

% --- Copyright page ---
\thispagestyle{empty}
\vspace*{\fill}
\begin{flushleft}
  \small
  \textcopyright\ 2026 Cédric Laubscher.\\[0.4em]
  This paper is made available under the
  \href{https://creativecommons.org/licenses/by/4.0/}{%
    Creative Commons Attribution 4.0 International (CC BY 4.0)} licence.\\[0.4em]
  You are free to share and adapt this text,
  provided you cite the author and the source.\\[1.2em]
  \textit{First edition, March 2026.}\\[0.4em]
  Typeset in \LaTeX\
\end{flushleft}
\vspace*{\fill}
\newpage

% --- Epigraph ---
\thispagestyle{empty}
\vspace*{4cm}
\epigraph{%
  Whoever we may be, from the atoms that form our molecules
  to our most intimate thoughts, we are nothing but
  the necessary expression of a simple and inexorable logic
  of coherence. A logic of causality that fulfils itself at
  every pulse. It is the Origin, the engine and the essence
  of the universe. The human being is merely its most vulnerable
  instrument.%
}{C.~Laubscher}
\newpage

% --- Table of contents ---
\frontmatter
\tableofcontents
\newpage

% ============================================================
\mainmatter
% ============================================================

% --- Foreword ---
\chapter*{Foreword \\ Why This Quest}
\addcontentsline{toc}{chapter}{Foreword}

There is a question I have never managed to put away in a drawer. Not the question of God, nor that of the meaning of life, although neither is far off. But a starker, more stubborn question: \emph{why does anything persist?}

\medskip

I am not talking about the cosmos in general, nor about life in particular. I am talking about the smallest and most ordinary thing there is: an electron. An electron exists. It has always existed, since the first seconds of the universe. It will still exist when our sun has gone out. It does not degrade. It does not age. It is rigorously identical to every other electron in the universe, at any place, at any moment. Why? Not ``how'' --- physics is perfectly capable of describing how an electron behaves. But \emph{why} is it there? Why this particular form, and not another? Why this stability, which seems to defy entropy\footnote{a measure of the disorder of a physical system; according to the second law of thermodynamics, it can only increase in an isolated system.}?

\medskip

This question led me to do something perhaps a little mad: to start from scratch. Not to correct existing physics. To suspend it entirely. Remove space. Remove time. Remove particles. And ask the question in its absolute nakedness: under these conditions, what \emph{can} exist?

\medskip

The answer I found --- or rather that I watched emerge, because one does not choose these things --- is at once simple and vertiginous. And it does not stop at the electron. It reaches all the way up to us.

\medskip

This paper is the account of that quest. Not a text. Not a scientific paper. An account: the way in which a simple idea, a pulse, led step by step to a vision of the world that I invite you to explore together.

\chapter*{Abstract}
\addcontentsline{toc}{chapter}{Abstract}

\textit{%
``Whoever we may be, from the atoms that form our molecules and our
cells to our most intimate thoughts, we are nothing but
the necessary expression of a simple and inexorable logic of coherence.
A logic of causality that fulfils itself at every pulse.
It is the Origin, the engine and the essence of the universe.
The human being is merely its most vulnerable instrument.''}

\medskip

This paper begins from a brutally simple question: \emph{why is there something rather than nothing?} Not as a poetic metaphor, but as a precise scientific problem.

\medskip

The Projective Dynamic Logo (PDL) is a fundamental research programme that attempts to answer it by starting from scratch: without presupposing the prior existence of space, time, or particles. Its starting point is an austere but powerful idea: for something to \emph{exist}, it is necessary and sufficient that it be able to repeat itself, that it form a stable cycle, a \emph{pulse}.

\medskip

The PDL shows that the smallest logically admissible pulse --- a structure of four points connected by six links --- is the archetype of the electron. The proton emerges as a richer hierarchical architecture, selected by purely combinatorial criteria. The great constants of physics --- $\hbar$, $\alpha$, $k_B$, $G$ --- are then no longer numbers given once and for all: they become coherence ratios, the signatures of the way in which each structure manages its own internal balance. It is important to be precise about the status of these derivations: $\alpha$ and $\hbar$ are obtained without any freely adjusted parameter; $G$, on the other hand, requires the calibration of a single parameter $\varepsilon$, fixed once only to reproduce its experimental value, then used consistently to simultaneously recover $\alpha$ and $k_B$. It is this mutual coherence --- one single parameter, three constants --- that constitutes the true test of the programme.

\medskip

No complex structure can close perfectly upon itself. There always remains a small \emph{coherence leakage}: an incompressible fraction of openness towards what lies beyond it. This leakage is the very mechanism by which gravitation emerges, and the structural condition of all transcendence. The human being --- organism, identity, thought --- does not escape this law: we are incomplete closures, kept alive by incessant exchanges with a world we do not entirely control. Our vulnerability is not an accident; it is the highest signature of the logic that governs the electron.

\medskip

Two recent results extend this picture. On the one hand, the PDL shows that a cavity built solely from finite signed graphs spontaneously reproduces the density-of-states law $\rho(\nu) \propto \nu^2$, the very law that underlies blackbody radiation --- which suggests that the dimensionality of space might itself be a consequence of coherence constraints. On the other hand, a structural sketch indicates that $(4,6)$ structures could provide a common language for general relativity and quantum mechanics; this result, clearly identified as preliminary, opens a path rather than a conclusion.

\medskip

This work is addressed to anyone who has ever wondered what it truly means to \emph{exist}, and why that question concerns physics as much as philosophy, the electron as much as oneself.

\newpage

% --- Introduction ---
\chapter*{Introduction \\ The Question Physics Dares Not Ask}
\addcontentsline{toc}{chapter}{Introduction}

There is a question that everyone has asked at least once, as a child or as an adult, gazing at the sky or staring into the void: \emph{why is there something rather than nothing?}

\medskip

Most of the time, one files this question away in the drawer of philosophy or religion, and moves on. Physics, for its part, seems to have decided not to concern itself with it: it describes \emph{how} things behave, not \emph{why} they exist. And yet something does not add up.

\medskip

Our two great theories of the world --- general relativity, which governs the cosmos and gravitation, and quantum mechanics, which governs the infinitely small --- are of a stupefying precision. They have been verified millions of times, and they are never truly wrong. But as soon as one pushes them to their limits --- at the heart of a black hole, at the moment of the Big Bang, at the edges of space --- their fundamental concepts collapse. Space no longer makes sense. Time no longer makes sense. Particles no longer make sense. This is not a problem of calculation; it is a problem of \emph{language}. Our most powerful conceptual tools simply do not know what they are talking about at these extremes.

\medskip

The Projective Dynamic Logo (PDL) takes this problem seriously. It does not propose yet another correction to the existing equations. It does something more radical: it \emph{suspends} everything. No space. No time. No particles given in advance. And it asks the bare question: under these conditions, what could still exist?

\medskip

The answer it proposes is at once simple and vertiginous. For something to truly exist --- not merely for an instant, but in a lasting, recognisable, real way --- it is necessary and sufficient that it form a cycle. A difference that returns to itself. A pulse. That is all.

\medskip

From this single principle, the PDL reconstructs step by step: the logical form of the electron, then that of the proton, then the great constants of physics. That handful of mysterious numbers which govern the entire universe and which nobody has ever truly been able to explain. It shows that these constants are not numbers fallen from the sky; they are \emph{coherence ratios}, the signatures of the way in which each structure manages its own internal balance.

\medskip

But the PDL does not stop at particle physics. It opens onto a broader vision: every thing that exists --- an electron, a proton, a human being, a thought --- is an incomplete closure. No structure can close perfectly upon itself. There always remains a leakage, an incompressible openness towards what lies beyond it. It is from this leakage that gravitation is born. It is this same openness that makes freedom, transcendence, and yes, our vulnerability possible.

\medskip

This paper is the first presentation of the PDL in English, written for the widest possible audience. You need no mathematical background to read it. Every technical concept will be translated into ordinary language before being used. Every analogy with human life will be offered as an invitation to reflect, never as a definitive certainty, but as an open proposition.

\medskip

The PDL is not a doctrine. It is a research programme: a rigorous and open way of asking questions that conventional physics sidesteps. Some of its results are solidly established. Others are precise conjectures, testable hypotheses. Others still are philosophical extrapolations, clearly identified as such. This honesty is, in our view, one of its greatest qualities.

\medskip

The PDL does not advance in isolation. It has engaged a formal dialogue with another fundamental research programme, the \emph{Theory of Objectivity} (TO), which seeks to identify the minimal logical conditions that any admissible universe must satisfy to allow the persistent existence of distinct objects. This dialogue --- recorded in a dedicated publication within the corpus --- has led to precise refinements: a more rigorous definition of active surfaces, a more explicit treatment of boundaries, and deeper reflection on what it means to ``exist for more than one observer''. These exchanges do not constitute a merger of the two programmes; they illustrate that the PDL, far from being an isolated edifice, is sufficiently open and structured to engage with other fundamental approaches without losing its internal coherence.

\medskip

We shall therefore begin at the beginning: nothingness, and the question of what can emerge from it. Not nothingness as a poetic backdrop --- nothingness as a logical challenge. What does it mean to exist, truly, when one is no longer permitted to presuppose anything whatsoever?

\medskip

The answer, we shall build together, step by step. And at the end of the path, perhaps, the phrase that guides this entire text will take on a new meaning:

\begin{quote}
\itshape Whoever we may be, from the atoms that form our molecules to our most intimate thoughts, we are nothing but the necessary expression of a simple and inexorable logic of coherence. A logic of causality that fulfils itself at every pulse. It is the Origin, the engine and the essence of the universe. The human being is merely its most vulnerable instrument.
\end{quote}

% ============================================================
\chapter{Nothing, Then Something}

Imagine that you had to erase everything. Not just the objects around you, not just the Earth or the Sun or the galaxies, everything. Space itself. Time itself. Nothing remains, not even an empty background on which something could appear. Under these conditions, what could possibly happen?

\medskip

Most physical theories dodge this question. They assume that there already exists a space-time, already some laws, already potential particles; then they describe how all that evolves. This is legitimate, and it is remarkably effective. But it cheats with the original question. The PDL refuses this evasion. It deliberately places itself \emph{upstream} of all that and poses the problem in all its nakedness: what is necessary and sufficient for there to be something rather than nothing, in a lasting way?

\subsection*{Instantaneous existence does not count}

First, a brutal observation: a thing that exists only for an instant does \emph{not really count}. If a difference appears and then disappears without leaving the slightest trace, with no possibility whatsoever of being recognised or retrieved, then it is logically indistinguishable from nothingness. What has existed only once, without any possibility of returning, is as if it had never taken place.

\medskip

This is not an abstract philosophical statement; it is a logical constraint. For something to be real in the strong sense of the term, it must be \emph{identifiable}, that is to say it must be retrievable, recognisable, re-instantiable. And this requires that it be capable of returning to itself.

\medskip

In the same way, a perfectly motionless state, absolutely uniform, without the slightest internal difference, is itself indistinguishable from nothingness: if strictly nothing happens, there is nothing to observe, nothing to distinguish, nothing to name.

\medskip

The only serious candidate for existence, in this austere framework, is therefore a difference that can repeat itself: a gap, a contrast, a distinction that is capable of returning to itself in a structured way. Not once; once is just chance. But a cycle, a return: that is something that can claim to exist.

\subsection*{The minimal pulse}

The PDL formalises this idea under the name of minimal binary pulse. Imagine two possible states: agreement and disagreement, yes and no, $+1$ and $-1$. A rule makes one pass from one to the other and back again: a perfect round trip, a cycle of length two. This is the absolute minimum of what one can call an event that lasts.

\medskip

Notice that, at this stage, there is still neither physical time, nor space, nor any measurable frequency. There is only the logical possibility of an alternation: something can return to what it was. That is all. But it is already infinitely more than nothing. I remember the moment when this idea struck me for the first time. Not in a laboratory, not in front of an equation, but in the silence of an ordinary night, wondering why an electron does not age. The answer I eventually glimpsed was so simple that it first seemed derisory. Then vertiginous.

\medskip

One can find an image in everyday life. A heart that beats does not exist because of the space in which it is located; it exists because it can repeat its cycle indefinitely. A melody does not exist because it occupies a place; it exists because one can recognise it, play it again, find it again. The PDL says: at bottom, everything that exists operates according to this same principle. Space and time come \emph{afterwards}, as tools for organising and comparing cycles, not before, as a pre-existing backdrop.

\subsection*{Why a single pulse is not enough}

An isolated pulse, however, poses a problem. If it is connected to nothing else, how can it be distinguished from another identical pulse? How can it \emph{refer} to itself? Without a relational context, it remains undefinable: two identical and isolated alternations are strictly indistinguishable from one another.

\medskip

For a distinction to be operational, it must be anchored in a network of relations, it must be able to situate itself with respect to something else. This is why the PDL introduces a second requirement: elementary entities must be linked to one another, and these relations must themselves form closed loops.

\medskip

The simplest closed loop is a triangle: three entities, each linked to the other two. A triangle is the smallest circuit capable of closing upon itself without appealing to anything external. And in a triangle, something remarkable can happen: the three relations can be mutually \emph{coherent} or \emph{incoherent}, depending on how their signs combine.

\medskip

The PDL then posits a central axiom: admissible structures are those that minimise triangular incoherence. A structure that fails to satisfy its own internal constraints cannot maintain itself; it dissolves. Only those configurations persist that manage to close upon themselves in a coherent way.

\subsection*{The first admissible structure}

We can now pose the concrete question: what is the smallest network of linked entities that simultaneously satisfies the two requirements — pulse and triangular coherence — without appealing to any external support whatsoever?

\medskip

The answer is unique, and its precision is surprising: four entities linked by six edges. Neither fewer, nor more. With three entities, one can form a coherent triangle, but one cannot propagate reference across a sufficiently rich network: a single loop is too fragile. With four entities and six edges, a logical tetrahedron, denoted (4,6), one reaches the first threshold at which everything fits together: all loops are coherent, the pulse cycle is possible, and the structure is self-sufficient.

\medskip

This (4,6) structure is not invented or chosen arbitrarily: it is \emph{derived} from the constraints. It is the first logically obligatory answer to the question: what can exist without presupposing anything else?

\medskip

In the next chapter, we shall see that this abstract structure (four points, six links, one pulse) is precisely what the PDL identifies as the logical archetype of the electron. Not because the electron \emph{is} a mathematical graph, but because the electron is, in nature, the first physical realisation of what the logic of coherence makes obligatory.

\medskip

Before we get there, let us pause for a moment on what we have just done. We have used neither space, nor time, nor energy, nor any law of physics. We have only asked: what can return to itself in a coherent way? And a precise structure has emerged, inevitable. That is the heart of the PDL: existence is not a backdrop. It is a problem of logical optimisation, and physical reality is its solution.

% ============================================================
\chapter{The First Pulse: The Electron}

We established in the previous chapter that the smallest structure capable of existing in an autonomous and coherent way is the (4,6) network: four entities, six links, one pulse cycle. This structure was not invented; it was derived from the most minimal constraints one can impose on anything that claims to exist durably.

\medskip

But one question arises immediately: does this purely logical construction have anything to do with the real world? The PDL’s answer is yes, and in a very precise way.

\subsection*{The electron, discreet star of the universe}

Among all known particles, the electron occupies a special place. It is strictly stable: unlike many particles that decay within fractions of a second, an isolated electron never transforms into anything else. It is universal: all the electrons in the universe are rigorously identical, at any place, at any moment. It is structurally simple: at all scales accessible to our instruments, it has no detectable substructure. And it is animated by an internal oscillation, the Compton frequency\footnote{a frequency related to the mass of the electron by $f = mc^2/h$; it is the natural rhythm at which the electron “beats” in its own rest frame.}, which is its own and is directly linked to its mass.

\medskip

These four characteristics — stability, universality, absence of internal structure, intrinsic oscillation — are precisely those one expects from the (4,6) structure. Stable because its internal coherence is maximal and does not degrade. Universal because it is the only solution to the logical problem posed. Without internal structure because it is minimal by construction. And animated by a pulse because that is its very definition.

\medskip

The PDL therefore proposes what it calls a minimal correspondence principle: the electron at rest is the first physical realisation of the logical closure (4,6). This is not a metaphysical identity. One does not say that the electron \emph{is} a graph. It is a precise and testable structural hypothesis.

\medskip

To understand it, one must know what is meant by a \emph{stationary regime}: a structure is said to be stationary when it beats, spins, oscillates, but without ever transforming into anything else. It changes state at each pulse, but it always returns to the same initial pattern. It is like a spinning top: it is in permanent motion, but its overall shape remains stable. By contrast, a \emph{non}-stationary structure degrades, disintegrates, or transforms into something else over time.

\medskip

In this framework, the question becomes: if something must exist in nature as the \emph{first} stationary regime, that is, the simplest structure that beats without ever degrading, what is the natural candidate? The PDL’s answer is the electron: stable, universal, without substructure, animated by a permanent internal oscillation. Everything one expects from the (4,6) closure, dressed in physics.

\subsection*{The cost of existing}

This correspondence is not just a pleasing analogy. It allows us to reread an equation everyone has seen without ever truly understanding it:

\begin{equation}
E = \hbar\,\omega
\end{equation}

Before reading it, a word on the term \emph{quantum} — plural: \emph{quanta}. In everyday life, energy seems able to take any value: one can heat water a little, a lot, very slightly. But at the scale of particles, nature does not work that way. Energy is not continuous like flowing water; it is transmitted in discrete packets, minimal portions that cannot be subdivided further. A quantum is precisely such a packet: the smallest amount of energy that a given structure can exchange in a single go. One cannot receive half a quantum, any more than one can receive half a coin in a transaction that demands it whole.

\medskip

In ordinary physics, the equation $E = \hbar\,\omega$ simply states that the energy of this packet is proportional to the frequency $\omega$ of the associated oscillation, with $\hbar$ (the reduced Planck constant, pronounced “h-bar”) as the proportionality factor. This is correct, but it is mute: why this relation? Where does $\hbar$ come from? Why this particular value?

\medskip

In the PDL, the same equation is read differently. $\omega$ is the pulse frequency of the (4,6) structure, the physical expression of its logical pulsation. And $\hbar$ is no longer a mysterious parameter imported from outside: it is the minimal cost of a coherence cycle, the amount of action required for the structure to loop back onto itself once. To exist, even in the simplest possible way, has a price. This price is paid at each pulse. And $\hbar$ is the currency.

\medskip

This rereading transforms what seemed to be an arbitrary law into a logical necessity. The Planck constant is not a curiosity of nature; it is the quantitative translation of the fact that no existence cycle is free.

\subsection*{Mass as a coherence budget}

If $\hbar$ is the cost of a cycle, then mass and energy can be reread as coherence budgets: measures of how many cycles a structure must maintain, and at what rate, in order to continue to exist.

\medskip

For the simple (4,6) structure, the electron, this budget is minimal and uniform. All its coherence is internal, well organised, without hierarchy. This is why the electron is so light: it pays the bare minimum to exist.

\medskip

Inertia, too, finds a new reading here. Why is it difficult to change the state of motion of a massive object? Because such an object is a structure whose internal coherence is dense and strongly organised. Modifying it requires reorganising a large number of cycles simultaneously, which demands a proportional investment of coherence. What we call mass is in reality the \emph{internal rigidity} of a logical closure.

\subsection*{A structure, not a point}

One of the most persistent habits of classical physics is to treat elementary particles as \emph{points}: entities without dimension, without interior, without structure. The PDL breaks with this habit. Not by assigning them a geometric size, but by recognising in them an internal relational organisation.

\medskip

The electron, in this framework, is not a point in space. It is a pattern of relations, four nodes, six links, a cycle that sustains itself. What matters is not where it is, but \emph{how it is organised}. The space in which it is located is an emergent description, convenient for comparing several electrons with one another, but no more fundamental than the relational structure that defines it.

\subsection*{The first instrument}

Let us return to the sentence that guides this paper. It says that we all are, from atoms to thoughts, instruments of a logic of coherence. The electron is the first of these instruments: the simplest way the universe has of concretely realising what logic makes obligatory.

\medskip

It is not there by chance. It is not there because an external law created it. It is there because, as soon as something must exist in a minimal, autonomous and coherent way, the (4,6) form is inevitable, and the electron is that form, dressed in physics.

\medskip

But the universe does not content itself with electrons. It builds far more complex structures: protons and neutrons, nuclei, atoms, molecules, stars, living beings. How does one pass from the minimal pulse to all this richness? That is the question of the next chapter.

\medskip

But the electron is not only defined by its structure; it is also defined by the way it responds when one interrogates it. And here again, the logic of coherence imposes a unique answer.

\subsection*{What measurement owes to coherence}

There is a question that quantum mechanics has been asking for a century without ever truly answering: why are quantum probabilities \emph{squares}? Why, when one measures the spin\footnote{an intrinsic property of quantum particles, analogous to a rotation on itself but with no classical equivalent; it takes only two discrete values, “up” or “down”.} of an electron, is the probability of obtaining a result proportional to the square of an amplitude, and not to the amplitude itself, nor to its cube?

\medskip

This rule, called Born’s rule, after the physicist Max Born who formulated it in 1926\footnote{%
    Formally, if a quantum system is in the state $|\psi\rangle = a\,|\!\uparrow\rangle + b\,|\!\downarrow\rangle$, where $a$ and $b$ are complex numbers satisfying $|a|^2 + |b|^2 = 1$, then the probability of obtaining the result “spin up” in a measurement is $|a|^2$, and that of obtaining “spin down” is $|b|^2$. It is the square of the modulus of the amplitude — and not the amplitude itself — that gives the probability. This rule applies to all observables in quantum mechanics, not just spin: the probability of finding a particle at a point in space, the probability that an atom emits a photon, the probability that a nuclear reaction occurs — all this is calculated via squared amplitudes. Why squares and not something else? In standard physics, this rule is simply postulated; the PDL proposes, for the first time in a discrete framework, a structural derivation of it.} 
is one of the cornerstones of all quantum physics. It is confirmed with extraordinary precision in millions of experiments. But in the standard framework, it is \emph{postulated}: one lays it down as a rule of the game, without explaining its origin. Why squares? Nobody truly knows. It is one of the questions that resisted me the longest. I circled around it for months before understanding that the answer was not in the existing equations, but upstream of them, in the very structure of what it means to measure something that beats.

\medskip

That is where the opening came from. The PDL gives an answer, and it is one of the programme’s most solid pieces. Within it, measuring the spin of an electron means coupling the (4,6) structure to a measuring device, itself modelled as a network of relations. This coupling is made through what are called mixed triangles: triplets of links that connect two nodes of the electron to one node of the device. These triangles can be coherent or incoherent, exactly like the internal triangles of the (4,6) structure. And the same principle applies: the system globally minimises incoherence.

\medskip

The device has two branches, say “spin up” and “spin down”. Each configuration of the global system is weighted according to the number of coherent mixed triangles it realises in each branch. In the limit where this optimisation is strong, the effective probabilities of ending up in one branch or the other converge towards a precise expression. And this expression is exactly Born’s rule: the probabilities are the squares of the amplitudes.

\medskip

This result was then consolidated by a Gleason-type proof\footnote{%
  Gleason’s theorem, established by the American mathematician Andrew Gleason in 1957, answers the following question: if one wants to assign probabilities to the results of quantum measurements in a coherent way — that is, in such a way that the probabilities do not depend on the measurement context chosen — what form can these probabilities take? The answer is remarkably constraining: in a Hilbert space of dimension at least three, the only possible assignment is precisely Born’s rule. Not one among others: the only one. This result is regarded as one of the deepest demonstrations in the foundations of quantum mechanics, because it shows that Born’s rule is not an arbitrary choice — it is the inevitable consequence of demanding minimal coherence in the way probabilities are assigned. The PDL establishes an analogous result in a discrete, combinatorial framework rather than in the continuous Hilbert-space (an abstract mathematical space in which quantum states are represented as vectors; the standard framework of quantum mechanics) of standard physics: this is what makes it a structural advance, and not a mere reformulation (\textit{Toward a Gleason-Type Uniqueness Theorem for Born's Rule in the PDL Framework}, C.~Laubscher, Zenodo, 2026).}.

\medskip

Concretely: under natural operational conditions —
phase invariance, rotational covariance, repeatability of measurement,
symmetry between the two branches — it is possible to show that Born’s rule
is the only assignment of probabilities compatible with the PDL structure.

\medskip

It is not merely that it emerges naturally. It is that no other
rule can exist in this framework without contradicting its own constraints.

\medskip

In other words: in a universe governed by the PDL’s logic of coherence, quantum probabilities \emph{could only be squares}. This is not a choice. It is a structural necessity.

\medskip

This result has far-reaching consequences. It means that the most fundamental rule of quantum mechanics — the rule that tells how the probabilistic world of the infinitely small translates into observable outcomes — is not an arbitrary convention. It is the inevitable consequence of a universe that optimises its relational coherence. The same logic that selects the electron as the first admissible structure also dictates the way it responds to a measurement.

% ============================================================
\chapter{Building Bigger: The Proton}

The electron is admirable in its simplicity. But the universe we inhabit does not reduce to isolated electrons. There are nuclei, atoms, molecules, and at the heart of all this there is the proton, about two thousand times heavier than the electron and with an internal complexity of a radically different kind.

\medskip

How does the PDL approach the proton? Not by postulating it, not by adding it in by hand. By \emph{deriving} it, by showing that, once one stacks (4,6) closures according to the same criteria of coherence and optimisation, a precise hierarchical architecture imposes itself.

\subsection*{One cannot just add electrons}

The first temptation would be to say: the proton is simply several (4,6) structures placed end to end. But that does not work. If each substructure seeks to maximise its own internal closure, the whole becomes too dense, too saturated. There is no longer room for a rich internal dynamics, no more flexibility, no more possibility of interaction with the outside. A sum of absolutes does not make a coherent whole.

\medskip

The PDL therefore imposes an additional rule: in a composite structure, there must be a hierarchy. Some regions are strongly closed; these are the stationary cores. Other regions are more open, more dynamic; this is the relational sea. This asymmetry between local closure and global openness is what allows the structure to remain coherent without freezing.\footnote{%
  Readers familiar with particle physics will recognise structural analogues here: the stationary cores correspond to \emph{valence quarks}, and the relational sea to the field of \emph{gluons} and \emph{virtual quarks} of quantum chromodynamics. But these technical terms are not necessary to follow the argument: the PDL reconstructs this architecture independently, without postulating it.}

\subsection*{Integers that are not parameters}

When the PDL applies its optimisation criteria to this architecture, namely: maximal coherence; minimal hierarchy; net charge compatible with that of the electron, it obtains a structure characterised by five precise integers:

\begin{itemize}
\item $n_u = 24$ and $n_d = 28$: the sizes of the two types of stationary cores (up and down quarks), expressed in number of elementary entities.
\item $r_{\text{val}} = 930$: the total number of internal links in the three valence cores.
\item $R_{\text{mer}} = 10\,087$: the number of links in the dynamic relational sea.
\item $r_{\text{total}} = 11\,017$: the total relational budget of the proton.
\end{itemize}

These numbers are not tuned to match experimental data. They are \emph{selected} by logical constraints: divisibility by four, controlled asymmetry between the cores, and maximisation of global coherence under a minimality constraint. The closeness to experimental results, in particular the $9\,\%$ stationary / $91\,\%$\footnote{The approximate 9\,\% stationary / 91\,\% dynamic split is consistent with the results of deep inelastic scattering (DIS) experiments. Combined measurements from the H1 and ZEUS experiments at the HERA collider have established that valence quarks carry only about 45\,\% of the proton’s total momentum, the rest being attributed to gluons and sea quarks. See: H1 and ZEUS Collaborations, \textit{Combined Measurement and QCD Analysis of the Inclusive $e^{\pm}p$ Scattering Cross Sections at HERA}, \textit{JHEP} 01 (2010) 109, arXiv:0911.0884 [hep-ex].}
dynamic partition of the proton’s energy, is an \emph{a posteriori} validation, not an adjustment.

\medskip

These five integers did not fall from the sky, nor were they adjusted by hand to fit the data. They are selected by the PDL optimisation functional under strictly combinatorial constraints\footnote{The derivation of these five integers, the systematic numerical simulations scanning the space of neighbouring configurations, and the local uniqueness conjecture are developed in full in \textit{On the Combinatorial Selection of the Proton Architecture in PDL --- Axiomatic Constraints, Integer Structures, and a Uniqueness Conjecture}, C.\,Laubscher, Zenodo, 2026. A formal proof of global uniqueness remains an open problem.}~: divisibility by four, maximal triangular coherence, stationary / dynamic fraction close to $9/91$, net charge $+1$, and simultaneous compatibility with the fine-structure constant. Systematic numerical simulations have been carried out\footnote{%
 These explorations are the subject of a dedicated article in the PDL corpus: \emph{On the Combinatorial Selection of the Proton Architecture in PDL — Axiomatic Constraints, Integer Structures, and a Uniqueness Conjecture} (C.~Laubscher, Zenodo, 2026). The simulations, performed on Google Colab, scanned the space of neighbouring configurations by varying the multiplicities $n_u$ and $n_d$ in steps of four and adjusting $R_\text{mer}$ in moderate intervals. In each case, the alternative configurations violated at least one constraint among the $9/91$ partition, the accuracy on $\alpha^{-1}$, or global triangular coherence. These results provide strong numerical support for the local uniqueness conjecture, without constituting a formal proof: a uniqueness theorem in the strict mathematical sense remains an open problem.}
to test the rigidity of this quintuplet. Changing a single integer, for example taking $r_\text{total}$ from $11\,017$ to $11\,018$, is enough to destroy the simultaneous agreement with the constants of nature. This is not a numerical coincidence: it is the signature of an extremely tight structural constraint. The proton architecture is therefore the most constrained candidate the programme has identified to date, but the formal and exhaustive demonstration of its absolute uniqueness remains an open problem and one of the most stimulating mathematical challenges that the PDL poses to specialists in combinatorial optimisation.

\medskip

I will not forget the moment when these five numbers ceased to be variables and became constraints. When one realises that one cannot change a single one without destroying everything, something changes in the way one looks at physics, and perhaps at reality as a whole.

\subsection*{The active surface and the golden ratio}

The proton does not live in isolation. It interacts with electrons, with other protons, and with radiation. For this it has an active surface: a subset of its relations projected outwards and made available for couplings with other structures.

\medskip

The solution to the optimisation problem for this surface brings out, in a conjectural but strongly constrained way, the golden ratio\footnote{a ratio that appears naturally in problems of self-similarity and optimal partition.} $\varphi = (1+\sqrt{5})/2$. Not as a mystical symbol, but as the value that optimises the partition between inside and outside in a hierarchical structure subject to the PDL constraints.

\subsection*{The constants of nature as coherence ratios}

It is at this stage that the PDL accomplishes something remarkable. The great physical constants here find a precise structural interpretation.

\medskip

The fine-structure constant $\alpha \approx 1/137$ is reinterpreted as the ratio between the active surface of the proton and its total relational budget: the fraction of coherence that the proton projects outward in order to interact with an electron.\footnote{%
  $\alpha$ (alpha) is one of the most famous constants in physics. It measures how strongly charged particles — such as the electron and the proton — attract or repel one another via the electromagnetic force. Its value, about $1/137$, is a dimensionless number: it does not depend on any system of units. Why precisely $1/137$ and not something else? Nobody knows within standard physics. This is precisely what the PDL seeks to explain.}
Boltzmann’s constant $k_B$ encodes the way coherence is distributed between microscopic and macroscopic levels of organisation.\footnote{%
  $k_B$ (pronounced “Boltzmann's constant”, after the Austrian physicist Ludwig Boltzmann, 1844–1906) is the bridge between two worlds: that of individual particles and that of the heat we feel. It translates what we call \emph{temperature} — a macroscopic notion, that of the thermometer — into the average agitation of particles at the microscopic scale. Without it, statistical physics and thermodynamics could not communicate.}
The gravitational constant $G$, finally, is obtained by hierarchical amplification of this same leakage rate $\varepsilon$ across eighteen levels of organisation\footnote{The detail of these eighteen filters, organised in three blocks of six levels (proton interior, proton-to-nucleus, nucleus-to-macroscopic matter), together with the derivation of the effective gravitational coupling, is developed in \textit{Coherence Leakage, Hierarchical Filtering, and the Exponent\,18 in the PDL Framework}, C.\,Laubscher, Zenodo, 2026.}. It is important to be precise about the status of this result: $G$ is not derived without a free parameter as $\alpha$ is. The value of $\varepsilon$ is, at this stage of the programme, adjusted to reproduce the experimental value of $G$; it is then this same calibrated $\varepsilon$ that is used to recover $\alpha$ and $k_B$ simultaneously. What the PDL establishes structurally is the \emph{form} of the link between $G$, $\alpha$ and the proton architecture; the numerical \emph{value} of $\varepsilon$ remains, for now, an empirical parameter. This distinction between a fixed logical architecture and partial numerical calibration is one of the marks of epistemic honesty of the programme.

\medskip

What is essential to retain: the laws of nature are not arbitrary rules inscribed in the cosmos. They are compact descriptions of the way a finite universe manages its coherence budgets at all scales, from the electron to the stars. They are the necessary conditions for a universe of partial closures to remain globally coherent.

\subsection*{An aside: what the PDL says about computation}

Readers familiar with computer science may have noticed a troubling resonance. The $(4,6)$ structure, a network of four nodes linked by six edges and subject to a discrete update rule, looks remarkably like a minimal cellular automaton. This is no accident. The PDL suggests that computation and physical existence share a deep common structure: in both cases, what runs stably is what \emph{exists}. A programme that loops indefinitely without halting is, in this framework, the software version of a closure; a programme that crashes is a structure that has been unable to maintain its internal coherence.

\medskip

This resonance is not just a metaphor. Recent work on quantum complexity and quantum circuits shows that the stability of a quantum computation depends precisely on the ability of logic gates to maintain entanglement (coherence) against environmental decoherence. The central problem of quantum computing — how to preserve coherence long enough to perform a useful computation — is, in the language of the PDL, the problem of maintaining an open closure without dissolving it. Quantum error-correcting codes, which detect and repair coherence violations without directly measuring the state of the system, can be read as technological implementations of the minimal-leakage principle.

\medskip

More generally, the PDL invites a reformulation of the fundamental question of computational complexity: why are some problems harder than others? In the relational framework of the PDL, the difficulty of a computation could be related to the coherence cost required to maintain all the constraints of the problem simultaneously. The complexity classes $\mathbf{P}$ and $\mathbf{NP}$ could then reflect a structural difference in the coherence budgets required, an open but precise conjecture.

\subsection*{A universe that builds itself}

The proton marks a decisive threshold. Beyond the first minimal closure that is the electron, here is a composite architecture that emerges from a hierarchy of organised pulses: stationary cores, a dynamic sea, an interaction surface, constants that are its imprints.

\medskip

But there is a limit to this nesting logic. No structure, however well organised, can close completely upon itself. There is always something that escapes, an incompressible leakage, an irreducible opening towards what lies beyond it. It is from this leakage, precisely, that gravitation is born. And it is this, too, that makes us vulnerable.

% ============================================================
\chapter{Nothing Ever Closes Completely: Leakage}

Up to this point, the PDL has built: first the minimal pulse, then the electron, then the proton. At each stage, the logic of coherence selects the most closed, the most stable, the most autonomous structure possible. One might think that the ultimate goal is perfect closure, a structure that would need nothing, that would maintain itself indefinitely without ever letting anything leak out.

\medskip

This perfect closure does not exist. And this is not a defect of the programme; it is one of its deepest results.

\subsection*{Irreducible incoherence}

Let us recall the principle of triangular coherence: in any admissible structure, triangles of relations must satisfy a condition of internal agreement. For the (4,6) structure of the electron, this condition can be satisfied perfectly: all triangles are coherent, no internal contradiction remains. This is why the electron is so stable.

\medskip

But as soon as one builds a composite structure, as soon as one assembles several elementary closures into a richer architecture, something changes. The number of triangular constraints to be satisfied simultaneously grows much faster than the number of available degrees of freedom. Beyond a certain level of complexity, it becomes \emph{mathematically impossible} to satisfy all constraints at once.

\medskip

Imagine a group of five friends who all want to get on perfectly with one another. Between two people, agreement is possible; one only has to find common ground. But between five people, there are ten pairs of relations to manage simultaneously. And experience shows that the larger the group becomes, the harder it is for everyone to get along with everyone else at the same time: some tensions become inevitable, not through ill will, but because the requirements of each relationship end up contradicting one another. One can minimise frictions, distribute them intelligently, but one cannot eliminate them all. This is exactly what the PDL says for physical structures: beyond a certain threshold of complexity, a residue of incoherence is structurally inevitable. This residue is coherence leakage.

\medskip

There is always a residue left: a fraction of triangles that cannot be coherent, whatever one does. The PDL quantifies this fraction by a parameter denoted $\varepsilon$, the coherence-leakage rate. For the proton, this rate is very small, but it is strictly positive. It cannot be brought to zero. It is structural.

\subsection*{What leakage becomes}

A first reaction would be to see $\varepsilon$ as a mere flaw, a residual imperfection without major consequence. The PDL says the opposite: this leakage is not lost. It is \emph{exported}.

\medskip

The coherence that cannot be absorbed inside a structure propagates outward, into the network of relations that surrounds that structure. It does not disappear; it becomes a constraint on neighbouring structures. On a large scale, this exported coherence manifests itself as a modification of what is called the metric, the way in which structures influence one another at a distance.

\medskip

In other words: gravitation is the macroscopic manifestation of coherence leakage\footnote{\textit{Resonance~--- Bekenstein, Hawking, and surface entropy (1972--1974). In 1972, Jacob Bekenstein showed that the entropy of a black hole is proportional to the area of its horizon~--- not its volume. Stephen Hawking completed this in 1974 by showing that black holes emit thermal radiation at a temperature that links simultaneously $G$, $\hbar$, and $k_B$~--- the three constants the PDL connects through a single leakage parameter~$\varepsilon$. This is not yet a derivation: the PDL does not yet explain Hawking radiation. It is a structural resonance: the idea that physical information resides on a boundary rather than in a volume is precisely what the proton's active surface instantiates in the PDL\@. A formal dialogue between these two frameworks remains to be built.}}. 
What Newton described as a force of attraction between masses, what Einstein reinterpreted as a curvature of space-time, the PDL reads as the collective propagation of the coherence irreducibly lost by each composite structure.\footnote{%
  This claim does not rest on intuition or analogy: it is the result of a detailed mathematical derivation, developed in several publications of the PDL programme (\textit{Coherence Leakage, Hierarchical Filtering, and the Exponent~18 in the PDL Framework}; \textit{A Structural Derivation of the Fine-Structure Constant from the PDL}; \textit{From Discrete Coherence Flux to Effective Fields}, all available on Zenodo). What gives it force, and what distinguishes it from a simple narrative reformulation, is the nature of the convergence obtained. The same formal framework, starting from the same logical constraints and without any \emph{ad hoc} fitted parameter, derives simultaneously and independently three quantities:
  \begin{itemize}
  \item the coherence-leakage rate $\varepsilon$, quantified at the scale of the proton;
  \item the fine-structure constant $\alpha \approx 1/137$, expressed as a combinatorial ratio on the proton structure;
  \item the gravitational constant $G$, obtained by hierarchical amplification of $\varepsilon$ across about eighteen levels of organisation.
  \end{itemize}
\textit{Epistemic note:} among these three quantities, $\alpha$ is derived without any freely adjusted parameter; the value of $\varepsilon$ is, by contrast, calibrated to reproduce $G$, and this same $\varepsilon$ is then used to test coherence with $\alpha$ and $k_B$. The strength of the result lies in this \emph{mutual coherence}: a single parameter, fixed once, simultaneously reproduces three constants that standard physics treats as entirely independent.
  These three values are compatible with experimental measurements. Their mutual coherence — the fact that they emerge from the same logical edifice without contradicting one another — constitutes the true keystone of the programme. In physics, recovering one fundamental constant by derivation is already remarkable. Recovering three simultaneously, from a single principle of coherence and without freedom of adjustment, is the most stringent test one can impose on a theory — and the PDL passes it.}

\subsection*{Hierarchical filtering and the exponent 18}

How can leakage so tiny at the scale of the proton produce an effect as observable as gravitation?

The PDL’s answer is a mechanism of \emph{hierarchical filtering}: leakage does not produce gravitation directly. It produces it after passing through eighteen successive levels of organisation. At each level, leakage from the lower level is slightly amplified and transmitted to the higher level. After eighteen steps, the signal, starting from almost nothing at the scale of a proton, has acquired the intensity we measure as gravitation.

Here, concretely, are the main stages of this cascade\footnote{%
  The detail of these eighteen filters, organised in three blocks of six, is developed in \emph{Coherence Leakage, Hierarchical Filtering, and the Exponent~18 in the PDL Framework} (C.~Laubscher, Zenodo, 2026).}:

\begin{itemize}
  \item Levels 1–6: the interior of the proton. Leakage is born at the heart of the proton architecture itself — between the valence cores, the dynamic sea, and the active surface. It is here that the parameter $\varepsilon$ is defined and quantified.
  \item Levels 7–12: from proton to atomic nucleus. Protons and neutrons assemble into nuclei. The leakage of each nucleon combines and is redistributed at nuclear level. Nuclear stability constraints constitute the filters of this block.
  \item Levels 13–18: from nucleus to macroscopic matter. Nuclei form atoms, atoms form molecules, molecules form macroscopic objects, then self-gravitating structures. At each step, accumulated leakage is amplified according to the laws of collective coherence until, at the end of the process, it produces an effective coupling of order $10^{-38}$ — that is, precisely what we measure for the gravitational constant $G$.
\end{itemize}

\medskip

This result deserves a pause. Conventional physics treats $\alpha$ and $G$ as two entirely independent constants\footnote{In 2025, a study published in \textit{AIP Advances} (Vopson \textit{et al.})
independently proposed that gravitation could emerge from an entropy-reduction process
in a computational universe --- a striking conceptual convergence with the PDL's
coherence-optimisation principle, reached by an entirely different route.}
: no one has ever known why they take the values they do, nor why they should have the slightest link between them. The PDL says the opposite: these two constants are two faces of a single combinatorial defect in the proton architecture. There is only one leakage, only one parameter $\varepsilon$; depending on whether one projects it towards the active surface (electromagnetic coupling) or towards hierarchical filtering (gravitation), one obtains $\alpha$ or $G$\footnote{Resonance~--- Verlinde and entropic gravity (2010) In 2010, Dutch physicist Erik Verlinde proposed that gravitation is not a fundamental force but emerges thermodynamically from entropy variations --- exactly as gas pressure emerges from the agitation of its molecules(E.\ Verlinde, \textit{On the Origin of Gravity and the Laws of Newton}, JHEP \textbf{04} (2011) 029, arXiv:1001.0785; M.\,M.\ Brouwer \textit{et al.}, \textit{First test of Verlinde's theory of emergent gravity using weak gravitational lensing measurements}, MNRAS 466, 2547 (2017)). Observational tests on more than 33{,}000 galaxies in 2016--2017 confirmed partial agreement with this prediction, without requiring dark matter. The PDL mechanism is structurally more precise: where Verlinde postulates an elastic volume entropy, the PDL proposes an explicit combinatorial derivation of the leakage $\varepsilon$, amplified across eighteen hierarchical levels. Both approaches converge towards the same central intuition --- gravitation as \emph{bookkeeping of information on boundaries} --- but by entirely independent paths.}.

\medskip

In other words, what is remarkable is not only the final result. It is that the same parameter $\varepsilon$, defined once and for all at the scale of the proton, simultaneously and without free parameter reproduces the value of $\alpha$ (the fine-structure constant) and the value of $G$ — two constants that standard physics treats as entirely independent of one another.

\subsection*{Biology as coherence optimisation}

If gravitation is the macroscopic manifestation of coherence leakage at the nuclear scale, one may ask what the PDL says about biological structures, which seem, by contrast, actively to \emph{resist} leakage.

\medskip

The living cell is, in this framework, a remarkable example of hierarchical closure. It maintains an electrochemical gradient across its membrane, selectively exchanges matter and energy with its environment, reproduces its internal organisation with extraordinary fidelity, and actively repairs its own incoherences. Contemporary molecular biology ceaselessly reveals the precision of these mechanisms: chaperone proteins, which help other proteins to find their correct conformation, are literally devices for repairing structural coherence. The mechanisms of DNA replication and correction, which reduce the copying error rate to less than one base in a billion, are biological implementations of the closure-optimisation principle.

\medskip

In this framework, life is not a mysterious phenomenon that escapes the laws of physics: it is the most elaborate form taken by coherence optimisation when closures acquire the capacity to reproduce themselves. Darwinian evolution, far from being a blind and arbitrary process, can be read as a search algorithm in the space of admissible closures: structures that maintain their coherence long enough to reproduce persist; the others dissolve. \emph{Natural selection} is, in this vocabulary, the filter that eliminates closures whose coherence budget is insufficient to ensure their own replication in a given environment.

\medskip

This Darwinian rereading does not weaken the theory of evolution — it anchors it in a more fundamental principle. Genetic variation, drift and selection continue to operate exactly as Darwin described them; what the PDL adds is an explanation of \emph{why} selection is possible, why some structures are more stable than others, and why biological complexity tends to increase: because a higher-level closure can manage its own coherence leakage more economically than a collection of independent closures.

\subsection*{Transcendence: a precise word}

The word \emph{transcendence} often evokes mystery or religion. The PDL proposes a rigorous and sober version of it.

\medskip

A structure is \emph{transcendent} in the PDL sense not because it belongs to another world, but because it cannot fully contain itself. Its coherence overflows. A part of what it is depends on conditions that lie beyond its own boundaries, in the network that surrounds it, in the structures with which it coexists, in a field of global coherence that no individual closure controls.

\medskip

In this sense, transcendence is not the exception; it is the rule. Everything that is sufficiently complex to be interesting is transcendent. Richness and openness go together. One cannot have both complexity and perfect closure at the same time.

\subsection*{Being open without dissolving}

This tension between closure and openness is universal in the PDL. Each structure must find a balance: closed enough to remain itself, open enough to interact, adapt, evolve. Too closed: it isolates itself and can no longer participate in the global dynamics. Too open: it loses its internal coherence and dissolves.

\medskip

We know this tension well in our daily lives. An individual who is too rigid ends up cutting themselves off from the world. An individual with no internal boundaries lets themselves be invaded and no longer knows who they are. Between the two there is a delicate space, that of living coherence, of the structure that maintains itself by adapting. The PDL says that this space is not a psychological metaphor: it is the structural condition of all complex existence, from the proton to the human being.

\subsection*{A universe that cannot close upon itself}

This is not a system that tends towards perfect equilibrium, towards total closure. It is a network of partially open structures that organise themselves with respect to one another, each exporting a fraction of its coherence into a global field that constrains them all.

\medskip

We are all, electrons, protons, human beings, nodes in this fabric. We export coherence. We receive it. We can neither do without it nor dissolve into it. It is this shared condition, this fundamental impossibility of total closure, that the next chapter will help us to reread in terms of causality and time.

% ============================================================
\chapter{Causality Revisited}

Why do things happen in a certain order? Why does a cause precede its effect? Why does time seem to go in only one direction? These questions have been at the heart of physics since Newton. And yet they remain strangely without a deep answer in current theories.

\medskip

Relativity says that time is relative to the observer. Quantum mechanics says that the future is not determined before measurement. But neither really explains \emph{why} there is a direction of time, why causes precede effects, why the world is not simply a frozen block of correlations without order. The PDL offers a different answer, and it begins by questioning the very definition of causality.

\subsection*{Causality as a chain: a convenient but limited image}

The classical image of causality is that of a chain: A causes B, B causes C, C causes D. Each link pushes the next. This image is intuitive, it works well to describe everyday phenomena, and it lies at the basis of all classical physics.

\medskip

But it becomes problematic as soon as one leaves simple situations. In quantum mechanics, an event does not have a single, localised cause: it results from a superposition of states. In general relativity, the temporal order between two events depends on the observer. And at the origins of the universe, the very notions of “before” and “after” lose their meaning.

\subsection*{Causality as global coherence}

In the PDL, causality is not a relation between two isolated events. It is a \emph{global} property: one event is causally linked to another not because it “pushes” it, but because both belong to the same coherence network, which cannot admit them simultaneously in just any order\footnote{{Resonance~--- ER\,=\,EPR: entanglement as geometry (2013--2024)} In 2013, Juan Maldacena and Leonard Susskind proposed the conjecture $\text{ER} = \text{EPR}$: two quantum-entangled particles would be connected by an Einstein-Rosen bridge (wormhole). In November~2024, a concrete and calculable realisation of this connection was published, explicitly deriving the bridge from entanglement in a conformal field theory. In the PDL, two coherent closures share \emph{mixed triangles} --- relations that span two distinct structures. This coupling is a discrete, combinatorial analogue of the $\text{ER} = \text{EPR}$ geometric connection: entanglement is not a mysterious state, it is a shared coherence between boundaries. Here again, this is a convergence, not a derivation.(J.\,Maldacena \& L.\,Susskind, \textit{Cool horizons for entangled black holes}, Fortschr.\ Phys.\ 61, 781 (2013), arXiv:1306.0533; concrete realisation: arXiv:2411.18485 (2024))}.

\medskip

Imagine a very full agenda: certain appointments cannot be moved without shifting others, not because they touch directly, but because they share common constraints — a room, a person, a deadline. It is not the first appointment that “forces” the second; it is the global structure of the agenda that imposes an admissible order. Remove the common constraint, and the order disappears with it. This is exactly what the PDL says: causality is not an arrow between two points, it is a global coherence constraint that is read off the structure of the entire network.

\medskip

What we call a \emph{cause} is a state that opens up a set of coherent transitions. What we call an \emph{effect} is the state that results from the transition that is most economical in terms of coherence, the one that best preserves the global integrity of the network. The causal order is therefore not an arrow imposed from outside; it is the signature of the way in which a network of closures manages its own maintenance.

\subsection*{Time as a reading of coherence}

Time, in this framework, is not a backdrop on which events unfold. It is an emergent reading of the order in which coherent configurations succeed one another. When a closure passes from one state to another while maintaining its coherence, it “moves forward in time”, not because it follows a universal calendar, but because it realises a sequence of admissible states in the only order possible given its internal constraints.

\medskip

This is not a negation of time; it is an explanation of its origin. Time is the way in which a universe of partial closures tells its own story of coherence maintenance. Where coherence can be preserved, time moves forward. Where it collapses — in singularities, at the origins — time no longer makes sense, precisely because there is no longer any admissible sequence.

\subsection*{Constants as causal signatures}

The fine-structure constant $\alpha$ does not just say “this is the strength of electromagnetism”; it says: “this is the fraction of coherence that a proton can share with an electron without either losing its integrity”. It is a causal constraint: if $\alpha$ were significantly different, atoms as we know them could not maintain their internal coherence, and chemistry, and life, would be impossible.

\medskip

The constants of nature are not arbitrary settings of one universe among other possibles; they are the necessary conditions for a universe of partial closures to remain globally coherent.

\subsection*{Determinism, chance, and freedom}

For simple structures — the electron, the proton — the sequence of admissible states is highly constrained. The dynamics looks like strict determinism. But for complex structures, those whose relational sea is vast, the number of admissible transitions explodes. There is no longer a single path; there are many, all coherent, all equally possible.

\medskip

What we call chance is the signature of this multiplicity. What we call freedom is the capacity, specific to the most complex structures, to \emph{represent} their own space of admissible transitions and to explore this space actively. Freedom and determinism are therefore not opposed in the PDL; they are two regimes of the same continuum, whose position depends on the internal richness of the structure considered.

\subsection*{Causality without a clockmaker}

What makes this picture particularly striking is the sobriety of the means it employs. The entire edifice of the PDL — constants, gravitation, time, causality — rests on only two ingredients: integers, which define the discrete structure of the relational network, and a geometric relation, the golden ratio $\phi$, which governs the redistribution of coherence between the core and the active surface of the proton. No adjustable continuous parameters. No constants introduced by hand. Two constraints, and an entire universe follows.

\medskip

The image that emerges from the PDL is that of a universe that does not need a clockmaker, an external law, a prime mover. Causality is not a rule imposed from outside; it is the inevitable consequence of the fact that coherent structures seek to maintain themselves. The order of the world is not programmed; it is \emph{negotiated}, at every moment, by millions of closures that simultaneously manage their own coherence and their inevitable leakage towards what lies beyond them.

\medskip

We are now ready to ask the question that this paper has only been preparing from the beginning: what does all this say about \emph{us}? What does it mean to be a human being in a universe governed by this logic of coherence?

\chapter{We, Vulnerable Instruments}

This chapter is the one that has cost me the most to write. Not for technical reasons, but because it is a matter of saying something honest about us, about what this logic implies if it is true. I wrote it, erased it, rewrote it. What you read here is the version in which I stopped censoring myself.

\medskip

We have followed the logic of coherence from nothingness to the proton, from the proton to the constants of nature, from the constants to causality and time. At each stage, the same law: to exist is to form a stable cycle, to manage a coherence budget, to accept irreducible leakage towards what exceeds us.

\medskip

It is time to ask the question this journey has prepared: what does this logic say about \emph{us}? Not about the electron, not about the proton, but about the human being, with their consciousness, emotions, choices, fragility.

\subsection*{We are high-level closures}

What follows is not a metaphor. It is the direct application of the same irreducible-leakage principle, mathematically demonstrated for the proton, quantified by $\varepsilon_G$, amplified across eighteen levels, to a composite structure of incomparably greater complexity: the human being.

\medskip

A human being is, from the PDL’s point of view, a composite structure of vertiginous complexity. Billions upon billions of electrons and protons, organised into atoms, molecules, cells, organs, systems. Each hierarchical level forms its own closures, manages its own coherence budgets, and exports its own leakage to higher levels and to the environment.

\medskip

But a human being is not only a stack of physical closures. They are also — and this is where the PDL touches something new — a closure that has developed the capacity to represent itself. A structure that can model its own functioning, anticipate its own future states, evaluate its own admissible transitions, and choose among them.

\medskip

It is this leap, from physical closure to reflexive closure, that distinguishes the human being from the electron or the proton. Not a break in the logic of coherence, but its most elaborate expression that we know: a structure that manages its coherence not only by reacting, but by \emph{anticipating}, by \emph{planning}, by \emph{meaning}.

\subsection*{Identity as closure maintained over time}

What is identity? In everyday life, it is what makes you \emph{you} today, tomorrow, in ten years, despite changes of body, ideas, relationships. In the language of the PDL, identity is the maintenance of a coherence pattern over time. Not the rigidity of a structure that never changes — the electron is rigid, but it has no identity in the human sense of the term. Rather, the persistence of an organisation that is reinstantiated at each cycle, slightly modified but recognisable, like a melody that can be played in different keys without ceasing to be itself.

\medskip

Maintaining one’s identity requires continuous effort. It is work of coherence: integrating new experiences without dissolving, adapting without betraying oneself, remaining open without losing the thread. This effort has a cost, a coherence cost in the literal sense of the PDL. And when this cost becomes too high, when contradictions accumulate, when losses exceed the capacity for integration, the closure wavers. What we call an identity crisis, a psychological collapse, or simply exhaustion, can be read as a deficit of internal coherence that the structure can no longer manage alone.

\subsection*{Memory, attention, effort}

Several dimensions of human experience find a precise echo in the vocabulary of the PDL. Not as reductions, but as \emph{resonances} that invite us to think differently.

\medskip

Memory can be read as the capacity to preserve structural regularities beyond a single cycle: coherence patterns that have been stabilised and can be reactivated. To remember is to reinstantiate a past closure in a present context. To forget is to let a pattern degrade to the point that it can no longer be retrieved.

\medskip

Attention looks like an allocation of coherence budget: directing one’s attention to something is to concentrate one’s internal coherence resources on a particular subnetwork, to the detriment of others. That is why attention is limited and fatiguable: it consumes something real.

\medskip

Effort, whether physical or mental, corresponds to the resistance one encounters when trying to modify a well-established closure — to change a habit, learn something new, overcome a fear. The more dense and old the closure, the more expensive it is in coherence to reorganise it. It is inertia, but applied to the mind: no longer the resistance of a physical mass, but the resistance of a tightly woven cognitive pattern.

\subsection*{Freedom: what a complex structure can do}

A structure as complex as the human brain has an internal space of coherent configurations so vast that no computation could enumerate it. Among these configurations, some are more stable, others more fragile. Some open up new possibilities, others close doors. Human freedom is the navigation in this space, guided by experience, constrained by biology and history, but real, because the space itself is real.

\medskip

This freedom is not infinite. It is bounded by the available coherence budget, by closures already established, by irreducible leakage towards an environment one does not control. But it is vast enough that what we call a choice, a deliberation, a decision, a commitment, is not an illusion. It is real navigation in a real space of admissible possibilities.

\subsection*{What medicine might read differently}

Contemporary medicine, especially in its most advanced branches, faces a paradox: the more precise it becomes in the molecular description of diseases, the more it struggles to grasp the patient as a coherent whole. The PDL offers here a complementary perspective, without claiming to replace medical biology.

\medskip

A disease, within the PDL, can be read as a disturbance of the coherence budget of a high-level closure. Cancer, for example, is a situation in which individual cells have “unlearned” how to respect the coherence constraints of the tissue to which they belong: they optimise their own local closure at the expense of the organism’s global coherence. It is precisely the situation described in the analogy of the group of friends: each member seeks to maximise their own relations to the detriment of collective coherence.

\medskip

Autoimmune diseases illustrate another regime: the immune closure confuses its own structures with foreign ones, generating internal triangular incoherence that turns against the organism. And neurodegenerative diseases such as Alzheimer’s can be read as a progressive degradation of the coherence pattern that constitutes identity: memory, interpreted as the capacity to reinstantiate past coherence patterns, dissolves as the underlying structures lose their integrity.

\medskip

These rereadings are not diagnoses; they are invitations to think differently about what medicine already does. Cell therapies, immunotherapies and precision medicine all seek, in their own way, to restore or strengthen coherence constraints that allow an organism to maintain itself. The PDL suggests that this medical project and the PDL’s fundamental physical project share a deep structure: in both cases, it is a matter of understanding how complex closures manage their coherence budget, and how to repair it when it is disturbed.

\subsection*{Vulnerability: the structural price of complexity}

The result here is structural, not metaphorical. In the PDL, coherence leakage $\varepsilon$ grows with hierarchical complexity: the richer a structure is, the larger its incompressible share of openness to the outside. The human being, the most complex composite structure we know, is the most extreme consequence of this.

\medskip

A human being is the most elaborate composite structure we know. Their coherence leakage is therefore immense, not in the physical sense of gravitational exponentiation, but in the existential sense: a considerable portion of what maintains a human being in coherence lies \emph{outside} them. In other people who recognise them. In the institutions that organise them. In the ecosystems that nourish them. In the languages that allow them to think. In the shared memories that give them a history.

\medskip

Remove these supports — total isolation, loss of recognition, destruction of the social fabric — and the human closure begins to degrade. This is not a metaphor: it is a well-documented clinical observation. Human vulnerability is not a contingent weakness that one could eliminate with more will or technique. It is the direct and inevitable consequence of our complexity. We are vulnerable because we are rich, because our coherence is so dense and so hierarchical that it cannot be maintained without a vast and diversified external support network.

\subsection*{Relation as a condition of existence}

If human coherence depends structurally on the outside, then relation is not a luxury or an ornament of life: it is a condition of existence.%
  \footnote{This statement has a precise technical content in the PDL: a structure whose coherence leakage exceeds its internal absorption capacity can only be maintained by importing coherence from its environment. Relational dependence is not a psychological property — it is a structural constraint on any high-complexity closure.}
To be in relation — with other people, with nature, with ideas, with a tradition — is to feed one’s own closure with external coherence. It is what allows the structure to maintain and develop itself.

\medskip

The PDL offers here, cautiously, a structural reading of what great ethical and spiritual traditions have always intuited: that the human being is fundamentally relational, that their dignity is inseparable from their fragility, and that to care for others is also to care for the coherence fabric that maintains us all in existence.

\subsection*{Instruments, not victims}

The guiding sentence of this paper calls the human being “the most vulnerable of instruments”. An instrument, in the PDL sense, is not passive: it is a structure through which a logic is realised, a form through which something universal finds a particular, local, irreplaceable expression.

\medskip

The electron is an instrument: through it, the logic of minimal coherence is realised for the first time in the physical world. The proton is an instrument: through it, that same logic builds its first complex architecture. And the human being is an instrument — the most elaborate, the richest, the most fragile — through which the logic of coherence reaches the level at which it can \emph{reflect upon itself}.

\medskip

This is, perhaps, what is most singular in the human condition: we are not only structures that maintain themselves, we are structures that \emph{know} that they maintain themselves, that can contemplate the logic that runs through them, that can choose to participate in it consciously. We are the place where the universe begins to understand itself.

\medskip

This does not place us above the logic of coherence; it plunges us more deeply into it. Our reflexivity is not an escape route: it is a responsibility. That of managing our coherence with care — our own, and that of the relational fabric on which we depend and which we help to weave.

\medskip

Let us take a concrete example. Someone loses a loved one. In ordinary language, we say they are “broken”. In the language of the PDL, one could say that their identity closure — that pattern of relations, memories, habits, projects, that constitutes their self — has undergone a massive disturbance. Part of the coherence network that defined them has disappeared with the other. They must rebuild themselves: find a new balance, redistribute their coherence budget, reorganise their internal relations. This is painful, not metaphorically but structurally: it is the real cost of a deep reorganisation of a complex closure. And what makes this reconstruction possible is precisely leakage: the incompressible openness towards others, towards the world, which prevents the closure from shutting itself in on its own suffering.

\medskip

Our vulnerability is not a flaw. It is the condition of our openness. And our openness is the condition of our existence.

\chapter*{Conclusion — Learning to Read the Logos}
\addcontentsline{toc}{chapter}{Conclusion}

We set out from nothingness. Not as a poetic backdrop, but as a logical challenge: what can exist without presupposing anything else? The PDL’s answer is radically sober: a difference that can repeat itself, a cycle that returns to itself, a pulse. That is all. And yet, from this single principle, an entire architecture has emerged: the electron, the proton, the constants of nature, gravitation, time, and even the human condition.

\subsection*{A programme that knows how to dialogue}

The PDL does not advance in isolation. It has engaged in a formal dialogue with another fundamental research programme: the Theory of Objectivity (TO), which seeks to identify the minimal logical conditions that any admissible universe must satisfy in order to allow the persistent existence of distinct objects for several observers\footnote{\textit{A Modal Dialogue Between the Projective Dynamic Logo and the Theory of Objectivity}, C.~Laubscher, Zenodo, 2026.}.

\medskip

This dialogue is not an academic curiosity. It has produced concrete refinements within the PDL programme: a more rigorous definition of active surfaces, a more explicit treatment of boundaries between a structure and its environment, and deeper reflection on what it means to “exist for more than one observer”. These refinements are not cosmetic; they have made certain claims of the programme more robust and more sharply delimited.

\medskip

What is worth retaining is that a research programme’s capacity to dialogue with its intellectual neighbours without dissolving — without losing its internal coherence — is itself a form of test. A programme that is too fragile collapses at the first critical glance. A programme that is too rigid closes itself off and no longer progresses. In this exchange, the PDL has shown that it can do what its own structures do: remain open without dissolving. Accept leakage towards what exceeds it, and emerge from it more precise.

\subsection*{What the PDL is not}

The PDL is not a finished theory that would replace current physics overnight. It is not a philosophical doctrine that claims to explain everything. It is a research programme: a set of precise mathematical constructions, testable hypotheses, conjectures honestly identified as such, and clearly posed open questions. This epistemic honesty is, in our view, one of its most precious qualities.

\subsection*{What the PDL changes}

It changes the way we think about existence: to exist is no longer a mysterious property, it is a performance, an active maintenance of coherence under constraint. It changes the way we think about the constants of nature: these numbers are no longer arbitrary; they encode the conditions under which a universe of partial closures can remain globally coherent. It changes the way we think about vulnerability: to be fragile is no longer a defect to be corrected, it is the structural condition of all complex existence. It changes, perhaps, the way we think about transcendence: no longer a separate beyond, but the structural impossibility of fully containing oneself.

\medskip

It also changes the way we think about what counts as a serious theory. A serious theory is not only a theory that predicts correct things; it is a theory that says precisely what would destroy it.

\subsection*{The PDL knows what could refute it. And it says so.}

First scenario: the constants $\alpha$ and $G$ are locked, in the PDL, by a precise proton architecture defined by five integers — $n_u = 24$, $n_d = 28$, $r_\text{val} = 930$, $R_\text{mer} = 10\,087$, $R_\text{tot} = 11\,017$. This architecture is unique: changing one of these integers, even by a single unit, destroys agreement with the measured values. If future measurements of $\alpha$ or $G$ were to deviate from the corridor defined by this architecture — not by a per cent, but by a minute fraction, for the window is extraordinarily narrow — the bridge between the two constants would collapse. This would not be a correctable detail. It would be a refutation.

\medskip

Second scenario: the PDL asserts that the $(24, 28, 930, 10\,087, 11\,017)$ architecture is the unique solution that simultaneously satisfies the programme’s constraints. If someone were to find an alternative architecture — different integers, a different logical proton, a different hierarchical closure — that reproduces the same constants with the same precision, then the claim to uniqueness would fall. The programme would remain interesting, but it would lose one of its strongest assertions.

\medskip

Two more recent results of the programme deserve mention, because they illustrate two very different ways in which the framework can be tested.

The first concerns light and cavities. In ordinary physics, a closed box — a cavity — can contain light. If one counts how many different vibration modes this light can adopt as a function of frequency, one obtains a well-known law: the number of modes grows like the square of the frequency, written $\rho(\nu) \propto \nu^2$. This law underlies blackbody radiation and, ultimately, all quantum field physics. The PDL has shown that a cavity built \emph{solely} from finite signed graphs, without presupposing the prior existence of space or of a wave equation, spontaneously supports periodic modes whose density follows exactly this same law, $\rho(\nu) \propto \nu^2$, in three dimensions\footnote{\textit{Discrete Cavity Modes and Logical Density of States in the PDL Framework}, C.~Laubscher, Zenodo, 2026.}. This result was not engineered to reproduce this law; it is its consequence. It is an independent, non-circular validation: the same framework that selects the electron and the proton reproduces, without additional parameter, a law that physics has regarded since Planck as a fundamental property of three-dimensional space. This suggests that the three-dimensional structure of space itself could emerge, in this programme, as a consequence of relational coherence constraints, rather than as a presupposed backdrop.

\medskip

The second result concerns one of the oldest challenges in theoretical physics: the unification of general relativity and quantum mechanics. These two theories describe the world with extraordinary precision, each in its own domain, but their fundamental languages are incompatible. The PDL sketches a framework in which this incompatibility could find a natural resolution: the $(4,6)$ structures that model the electron can be read as germs of Dirac spinors, while the emergent metric arising from coherence cycles provides an analogue of Einsteinian curvature\footnote{\textit{Towards an Einstein--Dirac Unification from the PDL Framework}, C.~Laubscher, Zenodo, 2026.}. Preliminary consistency checks — hydrogenic limit, gravitational redshift — have been carried out with encouraging results. It is essential to be clear about the status of this result: it is a structural sketch, not an accomplished unification. No physicist can yet claim that the PDL \emph{unifies} general relativity and quantum mechanics. What it proposes is a \emph{common language} in which the two theories can be formulated without principled contradiction, and that is already, if it withstands tests, a considerable advance.

\medskip

These two scenarios are not causes for fear; they are guardrails. They show that the PDL is not built to withstand everything. It is built either to be correct, or to be shown wrong in a precise, irrefragable way. That is the very definition of a serious scientific programme.

\subsection*{Energy, mechanical production, and the cost of closure}

The PDL has concrete implications for how we think about the production and transformation of energy. In the classical framework, energy is a conserved quantity that transforms from one form into another; thermodynamics describes its limits. In the PDL framework, energy is a coherence budget: the amount of internal organisation that a structure can maintain and transfer.

\medskip

This rereading has consequences for our understanding of machines. A heat engine, seen through the PDL, is a closure that exploits a coherence gradient between two reservoirs (hot and cold) to produce ordered work. Carnot efficiency, the maximal efficiency of a heat engine, is, in this vocabulary, the constraint on the fraction of coherence that a closure can transfer from a disorganised reservoir (hot) to a more organised closure (cold) without violating the second law of thermodynamics. What Boltzmann described as entropy — and what the PDL rereads as a measure of coherence leakage distributed over a very large number of degrees of freedom — is precisely the macroscopic translation of the PDL leakage parameter $\varepsilon$, aggregated over billions upon billions of elementary closures.

\medskip

Recent developments in energy — nuclear fusion, hydrogen fuel cells, superconducting materials — can all be read as attempts to build closures capable of maintaining or transferring coherence with minimal leakage. The challenge of thermonuclear fusion is precisely to maintain a plasma at more than one hundred million degrees in a sufficiently coherent state for fusion reactions to occur: it is a closure problem under the most extreme conditions. The recent advances of the ITER reactor and the promising results of inertial-confinement fusion (notably the net-positive energy records obtained at NIF in 2022–2023) illustrate concretely how fundamental the problem of plasma coherence management is, and how difficult it is to maintain an open closure without dissolving it.

\subsection*{The PDL in the face of recent discoveries}

A serious research programme does not live in a vacuum: it must confront the most recent experimental discoveries, not to \emph{explain} them prematurely, but to test whether its conceptual framework is compatible with them. Here are some of the most striking recent advances in physics and biology, read in the light of the PDL.

\begin{itemize}

\item \textbf{The images of the black holes M87* and Sgr A*.} In 2019 and 2022, the Event Horizon Telescope network produced the first direct images of an event horizon — that of the supermassive black hole in galaxy M87, then that at the centre of our own galaxy. These images show a luminous ring surrounding a shadow — exactly what Einstein’s general relativity predicts. In the PDL framework, a black hole is a region where the density of closures is so extreme that local coherence leakage can no longer propagate outwards: the event horizon\footnote{boundary beyond which no signal, not even light, can escape from a black hole.} is the boundary beyond which no coherence can be exported. The preliminary sketch of the Einstein–Dirac interface in the PDL — explicitly identified as speculative — thus finds a first geometric confrontation with observation.

\item \textbf{Gravitational waves and LIGO.} Since the first direct detection of gravitational waves by LIGO\footnote{Laser Interferometer Gravitational-Wave Observatory; an American detector that made the first direct observation of gravitational waves in 2015.} in 2015, more than a hundred events have been recorded — mergers of black holes and neutron-star pairs. These waves are, in classical physics, ripples in space-time curvature. In the PDL, they would be collective propagations of coherence leakage: perturbations of the global coherence field generated by extremely violent events involving very massive structures. The reinterpretation of $G$ as the result of hierarchical filtering across eighteen levels is precisely the kind of prediction that precision measurements of gravitational waves could, in the long run, test.

\item \textbf{Proton-mass measurements and the constants crisis.} Recent ultra-precise measurements of the proton mass (notably via Penning traps) and ongoing discussions about slight tensions between different measurements of the Hubble constant (the “Hubble tension”\footnote{a statistical disagreement between the two main methods of measuring the universe’s expansion rate; one based on the cosmic microwave background, the other on Cepheids and supernovae.}) raise questions about the stability of fundamental constants. If the PDL is right in tying $\alpha$ and $G$ to a precise proton architecture, then any temporal or spatial variation of these constants would be the signature of a variation in the architecture itself — which is difficult to conceive within the PDL, but nonetheless constitutes a valuable coherence test.

\item \textbf{Synthetic biology and the origins of life.} Advances in synthetic biology — in particular the laboratory creation of minimal cells capable of division (JCVI experiments since 2010, culminating in \textit{JCVI-syn3.0} in 2016 and \textit{syn3A} in 2021) — raise the question of what is minimally required for a biological structure to be “alive”. In the PDL framework, a minimal viable cell is the smallest biological closure capable of maintaining its internal coherence \emph{and} reproducing itself. The discovery that \textit{syn3.0} contains about 473 genes, of which the function of a third remains unknown, suggests that the space of minimal biological closures is not yet fully charted — and that the PDL could offer structural criteria for understanding why some configurations are viable and others are not.

\item \textbf{Non-equilibrium physics.} The work of Jeremy England (MIT) on \emph{adaptive dissipation}\footnote{Jeremy England (MIT, then GSK) has shown that any self-replicating system coupled to a thermal bath must necessarily produce a minimal amount of entropy, proportional to its growth rate and its internal irreversibility. This result establishes a formal link between thermodynamic dissipation and structural persistence, which the PDL reformulates in terms of relational coherence. See: J.~L.~England, \textit{Statistical Physics of Self-Replication}, \textit{J.~Chem.~Phys.}~139, 121923 (2013), doi:10.1063/1.4818538; and on adaptive dissipation: \textit{Nature Nanotechnology}~\textbf{10}, 919–923 (2015).} shows that physical systems subject to an external energy flux tend spontaneously towards configurations that maximise their dissipation — what he sometimes calls a “thermodynamics of evolution”. This idea is structurally close, though formulated in a different language, to the PDL principle that the most stable closures are those that best manage their coherence leakage to the environment. A formal dialogue between these two approaches remains to be built, but the conceptual convergence is striking.

\end{itemize}

\subsection*{Returning to the sentence}

Let us return, one last time, to the sentence that opened this text:

\begin{quote}
\itshape Whoever we may be, from the atoms that form our molecules to our most intimate thoughts, we are nothing but the necessary expression of a simple and inexorable logic of coherence. A logic of causality that fulfils itself at every pulse. It is the Origin, the engine and the essence of the universe. The human being is merely its most vulnerable instrument.
\end{quote}

This sentence is not a slogan. It is a working hypothesis, formulated precisely, anchored in rigorous mathematical constructions, open to refutation. But if it is correct, even partially, then it says something profound about what it means to exist. It means that every time you maintain a coherent thought, every time you keep a promise, every time you repair a damaged relationship, you participate in the same project as the electron that loops upon itself. The vocabulary changes, the scale changes, the richness changes, but the deep structure is the same: maintaining coherence, accepting leakage, remaining open without dissolving.

\subsection*{An invitation, not a doctrine}

The PDL is a young, ambitious, incomplete programme. It needs to be tested, criticised, refined — by physicists, philosophers, mathematicians, by all those who take an interest in it, who bring questions its authors have not yet thought to ask.

\medskip

If you are a physicist or mathematician, the technical publications, freely available on \href{https://zenodo.org}{Zenodo}, await you with their proofs, conjectures, and open questions. If you are a philosopher, you will find in them a precise and contestable ontology. And if you are simply someone who wonders what it means to live in this universe, then we hope this paper has given you a few new tools for posing that question more clearly.

\medskip

For perhaps that is what, in the end, it means to learn to read the Logos: not to decipher a sacred text, not to master a definitive theory, but to learn to recognise, in everything that endures, the silent work of coherence, and in every fragility, the signature of an irreducible opening towards what exceeds us.

\medskip

I do not know whether the PDL is correct. I know that it is coherent, that it is testable, and that it has not yet been found wanting. That is all one can ask of a research programme at this stage. What I hope is that it will at least have made you want to pose the question differently.

\chapter{If This Vision Is Correct\texorpdfstring{\,\ldots}{...}}

The PDL is a research programme, not a doctrine. But every serious programme must ask: \emph{if it is true, what does it change?} Here are some domains in which the logic of coherence, if it proves fundamental, invites us to rethink questions that were thought to be settled.

\subsection*{Biology and Darwin reread}

Darwinian evolution remains entirely intact in this framework. But it finds a deeper explanation in it: natural selection is not a blind process among others; it is the inevitable filter of a universe that preserves only closures capable of reproducing themselves. Variation, drift, selection operate exactly as Darwin described them; what the PDL adds is an answer to the question Darwin did not ask: why is there stability? Why do certain forms persist? Because the logic of coherence constrains them to.

\medskip

The living cell, in this framework, is the most sophisticated biological closure we know: it maintains its internal organisation, repairs its incoherences (DNA repair mechanisms reduce the copying error rate to less than one base in a billion), and reproduces itself. Life is not a mysterious phenomenon that escapes physics; it is the way in which the logic of coherence is realised at the molecular scale.

\subsection*{Medicine otherwise}

A disease, in this vocabulary, is a disturbance of an organism’s coherence budget. Cancer: cells that optimise their local closure at the expense of the tissue’s global coherence. Autoimmune diseases: an immune closure that confuses its own structures with foreign ones. Neurodegenerative diseases: a progressive degradation of the coherence pattern that constitutes identity and memory.

\medskip

These rereadings do not replace medical biology; they suggest that the medical project and the PDL project share a deep structure: in both cases, it is a matter of understanding how complex closures manage their coherence, and how to restore it when it is disturbed.

\subsection*{Quantum computing}

The central problem of quantum computing — maintaining the coherence of a quantum system long enough to perform a useful computation — is, in the language of the PDL, the problem of maintaining an open closure without dissolving it. Quantum error-correcting codes can be read as technological implementations of the minimal-leakage principle: they detect and repair coherence violations without destroying the system’s state.

\subsection*{Energy and machines}

A heat engine exploits a coherence gradient between two reservoirs to produce ordered work. The Carnot limit — the theoretical maximal efficiency of a heat engine — is, in this vocabulary, the constraint on the fraction of coherence that a closure can transfer without violating the second law. Today’s major energy challenges — nuclear fusion (for which the ITER reactor represents the most ambitious global effort), high-temperature superconductors, fuel cells — are all attempts to build closures capable of maintaining or transferring coherence with minimal leakage.

\subsection*{Recent observations}

The PDL does not yet predict precisely enough to be refuted by the most recent measurements, but several observational results resonate with its framework. The direct images of the black holes M87* and Sgr A* by the Event Horizon Telescope (2019 and 2022) show a precise geometric structure at the event horizon — precisely the boundary beyond which no coherence can be exported. Gravitational-wave detections by LIGO since 2015 confirm that gravitation propagates as a perturbation of space-time curvature — which the PDL rereads as large-scale propagation of coherence leakage\footnote{{Resonance~--- The island formula and the information paradox (2019--2025)} Since 2019, a series of results --- the \emph{island formula} --- has shown that the entropy of Hawking radiation follows a curve compatible with the preservation of quantum information (the Page curve), thanks to a surface contribution from \emph{inside} the black hole.\footnotemark\ This result, obtained in effective quantum gravity, suggests that information does not disappear into black holes: it is encoded on internal surfaces. The PDL offers a natural reading of this structure: the \emph{incompressible leakage} of a closure is never truly absorbed --- it is exported. If this reading is correct, the ``island'' of the formula could correspond to a domain where exported coherence remains traceable within the global relational network. This dialogue remains an open problem, identified here as a path rather than a conclusion.(G.\,Penington, \textit{Entanglement Wedge Reconstruction and the Information Paradox}, JHEP 09 (2020) 002; A.\,Almheiri \textit{et al.}, \textit{The entropy of Hawking radiation},
Rev.\ Mod.\ Phys.\ 93, 035002 (2021))}. And the current tensions in the measurement of the Hubble constant, what cosmologists call the \emph{Hubble tension}, could — \emph{without certainty} — signal that something in the deep structure of the constants is not yet understood.

\backmatter
\appendix

\chapter*{Appendix — Going Further}
\addcontentsline{toc}{chapter}{Appendix}

\subsection*{A. Glossary of Key Terms}

\begin{description}

\item[Pulse (binary pulsation)] Alternation between two distinct states that returns to itself in two steps. This is the minimal form of existence in the PDL: a difference capable of repeating itself.

\item[Coherence budget] The total amount of internal coherence that a structure can maintain simultaneously. Mass, energy, and certain physical constants are ways of measuring this budget at different scales.

\item[Closure] A relational structure that satisfies its own coherence constraints without appealing to an external support. To exist, in the PDL, is to instantiate a closure.

\item[Minimal stationary closure] The smallest structure that can both pulse and autonomously maintain its triangular coherence. In the PDL, this is the (4,6) structure, identified with the logical archetype of the electron at rest.

\item[Triangular coherence] Condition imposed on each triplet of mutually connected relations: their product must satisfy a precise sign rule. A coherent triangle is a closed loop without internal contradiction.

\item[Fine-structure constant ($\alpha$)] Measures the strength of electromagnetic coupling. In the PDL, it is reinterpreted as the ratio between the active surface of the proton and its total relational budget.

\item[Coherence leakage ($\varepsilon$)] The irreducible fraction of triangles that cannot be made coherent simultaneously in a composite structure. This leakage is exported to the environment and constitutes, on large scales, the basis of gravitational coupling.

\item[Hierarchical filtering] Mechanism by which tiny coherence leakage at the nuclear scale is amplified, level by level, until it produces an observable macroscopic effect — notably gravitation.

\item[Finite signed graph] A finite set of nodes linked by edges, each carrying a sign $+1$ or $-1$. This is the basic substrate on which all PDL closures are defined.

\item[Projective Dynamic Logo (PDL)] Fundamental research programme developed by Cédric Laubscher since 2024, which reconstructs physical reality from axioms defined on finite signed graphs, without presupposing space, time, or particles.

\item[Relational sea] In the proton architecture, the set of dynamic links that surround the stationary cores. It is analogous to the gluon and virtual-quark field of quantum chromodynamics.

\item[Coherence ratio] Way of reading fundamental physical constants in the PDL: not as arbitrary numbers, but as measures of how coherence is distributed between the interior, the surface, and the exterior of a given structure.

\item[(4,6) structure] Complete graph on four vertices and six edges: the first logically admissible closure under the PDL constraints. Identified with the archetype of the electron at rest.

\item[Active surface] Subset of the relations of a composite structure that is projected outward and available for couplings with other structures.

\item[Transcendence (in the PDL sense)] The structural property of any sufficiently complex closure: the impossibility of fully containing itself, the necessity of depending on a global coherence field that exceeds its own boundaries.

\end{description}

\subsection*{B. Map of Technical Publications}

All technical work on the PDL is freely available on \href{https://zenodo.org}{\textbf{Zenodo}}, under the name Cédric Laubscher. The following publications constitute the core of the programme:

\begin{enumerate}

\item \textit{The Emergence of Physical Reality within the Framework of the Projective Dynamic Logo (PDL)} — axioms, derivation of the (4,6) structure, correspondence principle with the electron.

\item \textit{Minimal Stationary Closures in the PDL Framework: Necessity of the (4,6) Block} — rigorous demonstration of the minimality of the (4,6) structure.

\item \textit{On the Combinatorial Selection of the Proton Architecture in PDL} — derivation of the integers 24, 28, 930, 10\,087, 11\,017 and uniqueness conjecture.

\item \textit{A Structural Derivation of the Fine-Structure Constant from the PDL} — combinatorial expression for $\alpha$ without an adjustable parameter.

\item \textit{Coherence Leakage, Hierarchical Filtering, and the Exponent 18 in the PDL Framework} — leakage mechanism and amplification to effective gravitational coupling.

\item \textit{From Discrete Coherence Flux to Effective Fields} — passage from discrete graphs to continuous field descriptions.

\item \textit{Discrete Cavity Modes and Logical Density of States in the PDL Framework} — standard density of states reproduced from finite graphs.

\item \textit{Logical Leakage as Self-Maintained Probability: A Topological Reformulation in the PDL Framework} — quantum probabilities reinterpreted as frequencies of coherence violations.

\item \textit{Towards a Schrödinger-Type Dynamics from the PDL} — link between coherence cycles and the Schrödinger equation.

\item \textit{Towards an Einstein--Dirac Interface in the PDL Framework} — bridges between the PDL, general relativity, and relativistic quantum mechanics.

\item \textit{PDL: Global Mapping of Structures, Results, and Open Problems} — overview of the programme, established results, conjectures, and open questions.

\item \textit{Whoever We May Be: Existence, Coherence, and Vulnerable Instruments} — ontological and philosophical synthesis of the PDL corpus for a mixed audience.

\end{enumerate}

% ============================================================
\end{document}